\documentclass[11pt]{article}
\usepackage{mathunal-base}
\mathunalfuente{serif}
\usepackage{amssymb}
\usepackage{graphicx}
\graphicspath{{fig/}}
\providecommand{\nota}[1]{\par\smallskip\noindent{\small\itshape\color{unalblue}#1}\par\smallskip}
\newcommand{\resp}{\underline{\hspace{3cm}}}
\newcommand{\vf}{\underline{\hspace{1.2cm}}}

\begin{document}

{\large\bfseries Segundo Parcial de Matemáticas Discretas --- 23 de mayo de 2022}

\medskip
\noindent Escuela de Matemáticas. Universidad Nacional de Colombia, Sede Medellín.

\medskip
\noindent\textit{Duración: 1 hora y 50 minutos. Puntaje (sólo para uso oficial): Ej.\ 1--2--3 (50), Ej.\ 4 (15), Ej.\ 5 (15), Ej.\ 6 (20). Solucione las preguntas en los espacios en blanco asignados a cada una. No está permitido sacar hojas en blanco ni ningún tipo de apuntes durante el examen. En los ejercicios 1 a 3 decida si las afirmaciones son verdaderas o falsas; en estos ejercicios no se tiene en cuenta el procedimiento. En los ejercicios 4 a 6 se tiene en cuenta el procedimiento.}

\nota{Transcripción: los grafos A y B se reproducen a partir de las figuras digitales del examen original.}

\begin{enumerate}
\item[]\textbf{Problema 1.} (15 puntos) Considere el siguiente grafo (Grafo A).
\begin{center}\includegraphics[height=4.2cm]{p2022_grafoA_0}\end{center}
Responda si las siguientes afirmaciones son verdaderas o falsas.
\begin{enumerate}[label=(\alph*)]
\item (4 puntos) El grafo A admite una caminata Euleriana. \hfill Respuesta: \vf
\item (4 puntos) El grafo A admite un circuito Euleriano. \hfill Respuesta: \vf
\item (4 puntos) El número cromático del grafo A es $c = 2$. \hfill Respuesta: \vf
\item (3 puntos) El grafo A es planar. \hfill Respuesta: \vf
\end{enumerate}

\item[]\textbf{Problema 2.} (20 puntos) Considere el grafo $G$ cuyos vértices son los números $\{2, 3, 4, 5, 6, 7, 8, 9, 10\}$. Dos vértices del grafo $G$ están unidos por un lado si y sólo si los correspondientes números son relativamente primos. Responda si las siguientes afirmaciones son verdaderas o falsas.
\begin{enumerate}[label=(\alph*)]
\item (5 puntos) El grafo $G$ admite una caminata Euleriana. \hfill Respuesta: \vf
\item (5 puntos) El grafo $G$ admite un circuito Euleriano. \hfill Respuesta: \vf
\item (5 puntos) El número cromático de $G$ es $c \geq 3$. \hfill Respuesta: \vf
\item (5 puntos) El grafo $G$ es planar. \hfill Respuesta: \vf
\end{enumerate}

\item[]\textbf{Problema 3.} (15 puntos) Considere el siguiente grafo (Grafo B).
\begin{center}\includegraphics[height=4.2cm]{p2022_grafoB_0}\end{center}
Responda si las siguientes afirmaciones son verdaderas o falsas.
\begin{enumerate}[label=(\alph*)]
\item (3 puntos) El grafo B admite una caminata Euleriana. \hfill Respuesta: \vf
\item (3 puntos) El grafo B admite un circuito Euleriano. \hfill Respuesta: \vf
\item (3 puntos) El número cromático del grafo B es $c \geq 3$. \hfill Respuesta: \vf
\item (3 puntos) El grafo B es planar. \hfill Respuesta: \vf
\item (3 puntos) Dados dos vértices distintos $v$ y $w$ del grafo B, existe una caminata Euleriana que empieza en $v$ y termina en $w$. \hfill Respuesta: \vf
\end{enumerate}

\item[]\textbf{Problema 4.} (15 puntos) Encuentre números enteros $s$ y $t$ tales que
\[
1 = 1247\,s + 245\,t.
\]
Respuesta: $s = \underline{\hspace{2cm}}$, \quad $t = \underline{\hspace{2cm}}$.

\item[]\textbf{Problema 5.} (15 puntos) Encuentre un número natural $u$ tal que
\[
377\,u \equiv 1 \pmod{1524}.
\]
Respuesta: \resp

\item[]\textbf{Problema 6.} (20 puntos) Alicia quiere recibir un mensaje seguro usando el protocolo RSA. Ella publica su llave pública $(N = 221,\ e = 5)$. Berta quiere mandarle el mensaje $M$ a Alicia. Para eso lo encripta usando la llave pública, y publica el mensaje encriptado:
\[
\text{Encriptado} = 124.
\]
El mensaje original $M$ es: \resp

\end{enumerate}

\vfill
\noindent{\small\itshape Transcripción digital --- MathUNAL}

\end{document}
